\documentclass[11pt,a4paper]{article}

% ── Packages ──────────────────────────────────────────────
\usepackage[utf8]{inputenc}
\usepackage[T1]{fontenc}
\usepackage{amsmath,amssymb,amsthm,mathtools}
\usepackage{hyperref}
\usepackage{cleveref}
\usepackage[margin=1in]{geometry}
\usepackage{enumitem}
\usepackage{booktabs}
\usepackage{natbib}
\bibliographystyle{abbrvnat}

% ── Theorem environments ──────────────────────────────────
\newtheorem{theorem}{Theorem}[section]
\newtheorem{proposition}[theorem]{Proposition}
\newtheorem{lemma}[theorem]{Lemma}
\newtheorem{corollary}[theorem]{Corollary}
\newtheorem{definition}[theorem]{Definition}
\newtheorem{remark}[theorem]{Remark}
\newtheorem{example}[theorem]{Example}

% ── Notation shortcuts ────────────────────────────────────
\newcommand{\R}{\mathbb{R}}
\newcommand{\E}{\mathbb{E}}
\newcommand{\Prob}{\mathbb{P}}
\newcommand{\Var}{\mathrm{Var}}
\newcommand{\Tr}{\mathrm{Tr}}
\newcommand{\dif}{\,\mathrm{d}}
\newcommand{\Veff}{V_{\mathrm{eff}}}

% ──────────────────────────────────────────────────────────
\title{When Does SGD Fail to Find Good Solutions?\\
Four Failure Modes and Necessary Conditions\\
for Implicit Regularization}

\author{C$\Lambda$ / Lightman\\
\texttt{Lightman.chang@gmail.com}}

\date{\today}

% ──────────────────────────────────────────────────────────
\begin{document}
\maketitle

% ══════════════════════════════════════════════════════════
\begin{abstract}
Stochastic gradient descent (SGD) is widely observed to find solutions that
generalize well, a phenomenon attributed to its implicit regularization through
gradient noise.  Yet the precise boundaries of this mechanism remain unclear:
\emph{when does SGD fail to find good solutions?}

We identify four explicit failure modes, each demonstrated by a concrete
counterexample with quantitative analysis.
\textbf{F1}: when the learning rate is too small, exponentially long mixing
times trap SGD in bad basins.
\textbf{F2}: when the learning rate is too large, a positive Lyapunov exponent
causes divergence.
\textbf{F3}: when the noise-to-curvature ratios $\Sigma/H$ are identical
across minima, SGD's selection mechanism degenerates.
\textbf{F4} (our main result): even when the dynamics mix and are stable, the
noise-driven preference for low gradient variance can systematically select
solutions that generalize \emph{poorly}---a noise-generalization reversal.
This occurs because good solutions that leverage diverse features exhibit
high per-sample gradient variance, while bad solutions that memorize exhibit
low variance.

From these four failure modes we derive four necessary conditions~(N1--N4)
that must hold simultaneously for SGD's implicit regularization to succeed.
The conjunction is not shown to be sufficient, and we classify this gap as
an open problem.  Our results delineate the fundamental limits of noise-driven
implicit regularization and highlight settings where explicit regularization
is indispensable.
\end{abstract}

% ══════════════════════════════════════════════════════════
\section{Introduction}
\label{sec:intro}

Modern deep learning relies heavily on stochastic gradient descent and its
variants.  A striking empirical observation is that SGD often finds solutions
that generalize well to unseen data, even in over-parameterized settings where
many interpolating solutions exist.  This phenomenon has been attributed to
SGD's \emph{implicit regularization}: the stochasticity inherent in
mini-batch sampling biases the iterates toward solutions with favorable
properties such as low gradient noise variance or low Jacobian complexity
\citep{smith2018dont,chaudhari2018stochastic,li2021happens}.

The mechanisms underlying this bias are increasingly well understood.  In the
non-interpolation regime, the multiplicative structure of gradient noise induces
an effective potential $\Veff$ that differs from the training loss, favoring
regions of low noise variance \citep{horsthemke1984noise,chaudhari2018stochastic}.
In the interpolation regime, solutions whose per-sample Hessians exceed a
critical threshold $2/\eta$ are destabilized by discrete dynamics
\citep{wu2023implicit}.

However, most analyses focus on \emph{when} implicit regularization succeeds,
leaving an equally important question open: \textbf{when does it fail?}
Understanding failure is essential for two reasons.  First, it identifies
settings where practitioners must employ explicit regularization.  Second, it
sharpens our theoretical understanding by delineating necessary conditions that
any successful account of SGD's generalization must satisfy.

\paragraph{Our contributions.}
We present four explicit failure modes~(F1--F4), each as a concrete
construction with quantitative analysis:
\begin{enumerate}[label=\textbf{F\arabic*},leftmargin=2em]
  \item \textbf{Insufficient mixing}: when $\eta$ is too small, Kramers escape
        times grow exponentially, trapping SGD in bad basins for any
        feasible training duration (\Cref{prop:f1}).
  \item \textbf{Discrete instability}: when $\eta$ is too large, the top
        Lyapunov exponent becomes positive and SGD diverges (\Cref{prop:f2}).
  \item \textbf{Degeneracy of selection}: when $\Sigma/H$ is identical at all
        minima, SGD's noise-driven preference provides no selective advantage
        (\Cref{prop:f3}).
  \item \textbf{Noise-generalization reversal}: the good solution has high
        gradient noise (due to feature diversity), causing SGD to
        \emph{systematically} converge to the bad solution (\Cref{thm:f4}).
\end{enumerate}
Failure mode~F4 is our main result.  It reveals a fundamental limit: SGD's
preference for low gradient variance can conflict with generalization when
good solutions use diverse per-sample features and bad solutions memorize
uniformly.

From these four failure modes we extract four necessary conditions~(N1--N4)
for implicit regularization (\Cref{thm:necessary}).  We emphasize that their
conjunction is not proven sufficient, and we classify the sufficiency question
as an open problem.

\paragraph{Paper structure.}
\Cref{sec:prelim} introduces notation, the SGD setup, and background on noise
selection in one dimension (as known results).
\Cref{sec:failure} presents the four failure modes with full constructions and
proofs.
\Cref{sec:necessary} states the necessary conditions framework.
\Cref{sec:reparam} establishes reparameterization invariance of $\Sigma/H$
and its implications.
\Cref{sec:discussion} discusses practical implications, connections to prior
work, and open questions.
\Cref{sec:conclusion} concludes.

% ══════════════════════════════════════════════════════════
\section{Preliminaries}
\label{sec:prelim}

% ── 2.1 ──────────────────────────────────────────────────
\subsection{SGD Setup}
\label{sec:sgd-setup}

Let $\Theta \subseteq \R^d$ be the parameter space and let
$L(\theta) = \frac{1}{n}\sum_{i=1}^{n} \ell_i(\theta)$ be the empirical risk,
where each $\ell_i$ is the loss on the $i$-th training example.
SGD with learning rate $\eta > 0$ iterates
\begin{equation}\label{eq:sgd}
  \theta_{t+1} = \theta_t - \eta\, \nabla \ell_{z_t}(\theta_t),
  \qquad z_t \sim \mathrm{Uniform}(\{1,\dots,n\}) \text{ i.i.d.}
\end{equation}

% ── 2.2 ──────────────────────────────────────────────────
\subsection{Gradient Noise Variance}
\label{sec:noise-variance}

The gradient noise variance at $\theta$ is defined as
\begin{equation}\label{eq:sigma}
  \Sigma(\theta)
  = \frac{1}{n}\sum_{i=1}^{n}\bigl[\ell_i'(\theta)\bigr]^2
    - \bigl[L'(\theta)\bigr]^2.
\end{equation}
At a local minimum $\theta^*$ where $L'(\theta^*)=0$, this simplifies to
$\Sigma(\theta^*) = \frac{1}{n}\sum_{i=1}^{n}[\ell_i'(\theta^*)]^2$.
In the multi-dimensional setting, $\Sigma(\theta)$ denotes the covariance
matrix of the stochastic gradient; in one dimension it reduces to a
non-negative scalar.

% ── 2.3 ──────────────────────────────────────────────────
\subsection{Convergence and Good Solutions}
\label{sec:convergence}

We distinguish three levels of convergence:
\begin{enumerate}[label=(C\arabic*)]
  \item \emph{Optimization convergence}: $L(\theta_t) \to L^*$ almost surely.
  \item \emph{Distributional convergence}: $\theta_t \xrightarrow{d} \mu_\eta$
        for some stationary measure $\mu_\eta$.
  \item \emph{Absorption convergence}: $\theta_t \to \theta^*$ almost surely
        for some local minimum $\theta^*$.
\end{enumerate}

A solution $\theta^*$ is considered ``good'' if it satisfies one or more of
the following criteria:
\begin{enumerate}[label=(G\arabic*)]
  \item \emph{Low noise ratio}: $\Sigma(\theta^*)/H(\theta^*)$ is small
        relative to other local minima, where $H(\theta^*) = L''(\theta^*)$.
  \item \emph{Low Jacobian complexity}: in the interpolation regime,
        $\mathcal{C}(\theta^*) = \|J\|_F / \sqrt{n}$ is small.
  \item \emph{Leave-one-out stability}: the stability coefficient $\beta$ is
        small.
  \item \emph{Generalization}: the gap $L_{\mathrm{test}} - L_{\mathrm{train}}$
        is small.
\end{enumerate}
A central theme of this paper is that SGD's noise-driven bias favors~(G1),
and that the chain (G1)$\Rightarrow$(G3)$\Rightarrow$(G4) holds under mild
regularity conditions---but that (G1) can \emph{conflict} with~(G4) in
adversarial constructions (failure mode~F4).

% ── 2.4 ──────────────────────────────────────────────────
\subsection{Background: SGD Noise Selection in One Dimension}
\label{sec:background-1d}

We review known results on the stationary distribution of SGD in one
dimension, which form the analytical foundation for our counterexamples.
These results are not new; we state them for completeness and to fix notation.

\begin{proposition}[FPE steady state; \citealt{horsthemke1984noise,chaudhari2018stochastic}]
\label{prop:fpe}
Consider the one-dimensional SDE approximation of SGD with temperature
$\tau = \eta$:
\[
  d\theta = -L'(\theta)\dif t + \sqrt{\tau\,\Sigma(\theta)}\dif W_t.
\]
Assume $L,\ell_i \in C^2(\R)$, $L$ is coercive ($L(\theta)\to+\infty$ as
$|\theta|\to\infty$), and $\Sigma(\theta)$ is strictly positive and bounded:
$0 < \sigma_{\min}^2 \le \Sigma(\theta) \le \sigma_{\max}^2$.
Then the Fokker--Planck equation
\[
  \frac{\partial p}{\partial t}
  = \frac{\partial}{\partial\theta}\bigl[L'(\theta)\,p\bigr]
    + \frac{\tau}{2}\frac{\partial^2}{\partial\theta^2}\bigl[\Sigma(\theta)\,p\bigr]
\]
has a unique stationary solution
\begin{equation}\label{eq:fpe-steady}
  p(\theta)
  = \frac{1}{Z}\cdot\frac{1}{\Sigma(\theta)}
    \cdot\exp\!\Bigl(-\frac{2}{\tau}\int^{\theta}
      \frac{L'(s)}{\Sigma(s)}\dif s\Bigr),
\end{equation}
or equivalently $p(\theta) \propto \exp\!\bigl(-\frac{2}{\tau}\Veff(\theta)\bigr)$,
where the effective potential is
\begin{equation}\label{eq:veff}
  \Veff(\theta) = \int^{\theta}\frac{L'(s)}{\Sigma(s)}\dif s
    + \frac{\tau}{2}\ln\Sigma(\theta).
\end{equation}
\end{proposition}

\begin{proof}[Proof sketch]
Setting $\partial p/\partial t = 0$ in the Fokker--Planck equation, define the
probability current $\mathcal{J} = -L'p - \frac{\tau}{2}(\Sigma p)'$.
Since $\frac{d\mathcal{J}}{d\theta} = 0$, the current $\mathcal{J}$ is
constant.  Coercivity of $L$ and boundedness of $\Sigma$ force $p(\theta)\to 0$
as $|\theta|\to\infty$, which implies $\mathcal{J}\equiv 0$.
Expanding $\mathcal{J}=0$ gives
\[
  L'p + \frac{\tau}{2}\bigl[\Sigma' p + \Sigma p'\bigr] = 0
  \quad\Longrightarrow\quad
  \frac{p'}{p} = -\frac{2L'}{\tau\Sigma} - \frac{\Sigma'}{\Sigma}.
\]
Recognizing the right-hand side as
$-\frac{d}{d\theta}\bigl[\frac{2}{\tau}\int^{\theta}\frac{L'}{{\Sigma}}\dif s
+ \ln\Sigma\bigr]$, we integrate to obtain
$\ln p = -\frac{2}{\tau}\int^{\theta}\frac{L'}{\Sigma}\dif s - \ln\Sigma + C$.
Exponentiating yields~\eqref{eq:fpe-steady}.  The normalization constant
$Z<\infty$ is guaranteed by coercivity and boundedness of~$\Sigma$.
\end{proof}

The key feature of~\eqref{eq:fpe-steady} is the prefactor $1/\Sigma(\theta)$:
the stationary distribution assigns higher density to regions of \emph{low}
gradient noise variance.  This is the mathematical basis of SGD's noise-driven
selection.

\begin{proposition}[Generalization gap via leave-one-out;
  known result]
\label{prop:gap}
At a local minimum $\theta^*$ of $L$ with $H=L''(\theta^*)>0$, the
leave-one-out generalization gap satisfies
\begin{equation}\label{eq:gap}
  \mathrm{Gap}(\theta^*) \approx \frac{\Sigma(\theta^*)}{H(\theta^*)(n-1)}.
\end{equation}
\end{proposition}

\begin{proof}[Proof sketch]
Removing the $j$-th sample gives $L'_{-j}(\theta^*) = -\ell_j'(\theta^*)/(n-1)$
and $L''_{-j}(\theta^*) = (nH - \ell_j''(\theta^*))/(n-1) \equiv H_{-j}$.
By the implicit function theorem,
$\theta^*_{-j} - \theta^* \approx \ell_j'(\theta^*)/((n-1)H_{-j})$.
The leave-one-out loss difference is
$\ell_j(\theta^*_{-j}) - \ell_j(\theta^*)
 \approx [\ell_j'(\theta^*)]^2/((n-1)H_{-j})$.
Averaging over $j$ and using $L'(\theta^*)=0$ to identify
$\frac{1}{n}\sum_j[\ell_j'(\theta^*)]^2 = \Sigma(\theta^*)$, we obtain
$\mathrm{Gap} \approx \Sigma(\theta^*)/((n-1)H(\theta^*))$.
\end{proof}

% ══════════════════════════════════════════════════════════
\section{Four Failure Modes}
\label{sec:failure}

We now present the four failure modes.  Each is an explicit construction with
quantitative analysis, demonstrating a specific mechanism by which SGD's
implicit regularization breaks down.  We label
the violations of necessary conditions to be derived in \Cref{sec:necessary}.

% ── 3.1 F1 ───────────────────────────────────────────────
\subsection{F1: Insufficient Mixing}
\label{sec:f1}

\begin{proposition}[Failure by insufficient mixing]
\label{prop:f1}
There exists a one-dimensional loss landscape with two local minima---one good
and one bad---such that for $\eta = 0.01$, the expected time for SGD to escape
the bad basin exceeds $10^{1087}$ iterations.
\end{proposition}

\paragraph{Construction.}
Consider a symmetric double-well potential with barrier height $\Delta V = 100$
separating two wells:
\begin{itemize}[nosep]
  \item Well~A (good): Hessian $H_A = 2$, noise variance $\Sigma_A = 1$.
  \item Well~B (bad): Hessian $H_B = 8$, noise variance $\Sigma_B = 8$.
\end{itemize}
The noise ratios satisfy $\Sigma_A/H_A = 0.5 < \Sigma_B/H_B = 1$, so
the effective potential~\eqref{eq:veff} favors well~A (the good solution).
However, the barrier is so high that SGD cannot reach the steady state in
any feasible number of iterations.
SGD is initialized in well~B.

\begin{proof}
We apply the Kramers escape formula for the SDE with multiplicative noise
\citep{kramers1940brownian}.  The effective barrier for escape from well~B
to well~A is computed from the effective potential~\eqref{eq:veff}.

Near well~B, $\Veff$ has a local minimum.  The barrier in the effective
potential from~B to the saddle point is
\[
  \Delta U_{B\to\mathrm{saddle}}
  = \frac{\Delta V}{\Sigma_B} = \frac{100}{8} = 12.5.
\]
The Kramers escape time from well~B is
\[
  \tau_{B\to A}
  \sim \frac{1}{\eta}\exp\!\Bigl(\frac{2\,\Delta U_{B\to\mathrm{saddle}}}{\eta}\Bigr)
  = \frac{1}{\eta}\exp\!\Bigl(\frac{25}{\eta}\Bigr).
\]
Setting $\eta = 0.01$:
\[
  \tau_{B\to A}
  \sim 100\cdot\exp(2500)
  = 100\cdot 10^{2500/\ln 10}
  \approx 10^{1087}.
\]
This far exceeds any feasible training budget.

The key point is structural: regardless of the precise constant in the
exponent, for any barrier $\Delta > 0$, the Kramers time scales as
$\exp(\Theta(1/\eta))$, which grows super-exponentially as $\eta\to 0$.
With $\eta = 0.01$, the SGD iterates remain trapped in well~B for any
practical number of iterations, despite well~A being the preferred basin
under the stationary distribution.

This failure violates the mixing condition: the total training time $\eta T$
is negligible compared to the mixing time $\tau_{\mathrm{mix}}$.
\end{proof}

\paragraph{Which condition is violated.}
Condition~N1 ($\eta T \gg \tau_{\mathrm{mix}}$) is violated.
SGD has not reached its stationary distribution and remains trapped in the
initialization basin.

% ── 3.2 F2 ───────────────────────────────────────────────
\subsection{F2: Discrete Instability}
\label{sec:f2}

\begin{proposition}[Failure by discrete instability]
\label{prop:f2}
There exists a quadratic loss with two per-sample Hessians such that
$\eta = 0.15$ yields a positive top Lyapunov exponent, causing SGD to diverge.
\end{proposition}

\paragraph{Construction.}
Let $n=2$ with per-sample losses $\ell_1(\theta)=\frac{1}{2}\cdot 18\cdot\theta^2$
and $\ell_2(\theta)=\frac{1}{2}\cdot 2\cdot\theta^2$, so that
$L(\theta) = \frac{1}{2}\cdot 10\cdot\theta^2$.
The per-sample Hessians are $h_1 = 18$ and $h_2 = 2$.
The stability threshold is $\eta_c = 2/h_{\max} = 2/18 \approx 0.111$.
We set $\eta = 0.15 > \eta_c$.

\begin{proof}
The SGD iteration at the origin is $\theta_{t+1} = (1 - \eta\,h_{z_t})\,\theta_t$,
where $z_t\in\{1,2\}$ with equal probability.  The multiplicative factors are
\[
  1 - \eta\, h_1 = 1 - 0.15\times 18 = -1.7,
  \qquad
  1 - \eta\, h_2 = 1 - 0.15\times 2 = 0.7.
\]
Taking absolute values and logarithms:
\[
  \ln|1-\eta\,h_1| = \ln 1.7,
  \qquad
  \ln|1-\eta\,h_2| = \ln 0.7.
\]
The top Lyapunov exponent is the expected log-contraction rate:
\[
  \Lambda
  = \frac{1}{2}\bigl(\ln 1.7 + \ln 0.7\bigr)
  = \frac{1}{2}\ln(1.7\times 0.7)
  = \frac{1}{2}\ln 1.19
  \approx \frac{1}{2}\times 0.1740
  = 0.0870.
\]
Since $\Lambda > 0$, the strong law of large numbers gives
\[
  \frac{1}{T}\ln|\theta_T|
  = \frac{1}{T}\ln|\theta_0| + \frac{1}{T}\sum_{t=0}^{T-1}\ln|1-\eta\,h_{z_t}|
  \;\xrightarrow{a.s.}\; \Lambda > 0.
\]
Therefore $|\theta_t| \sim e^{0.087\,t} \to \infty$ almost surely: SGD diverges.
\end{proof}

\paragraph{Which condition is violated.}
Condition~N2 ($\eta < 2/\lambda_{\max}$) is violated.
Here $\lambda_{\max} = h_1 = 18$ and $\eta = 0.15 > 2/18$.

% ── 3.3 F3 ───────────────────────────────────────────────
\subsection{F3: Degeneracy of Selection}
\label{sec:f3}

\begin{proposition}[Failure by degenerate selection]
\label{prop:f3}
There exist two local minima, one good and one bad, with
$\Sigma_A/H_A = \Sigma_B/H_B$, such that the SGD stationary distribution
assigns equal probability to both: $\pi(A)\approx \pi(B) \approx 1/2$.
\end{proposition}

\paragraph{Construction.}
Consider a symmetric double-well loss (symmetric about the midpoint between
wells~A and~B) with:
\begin{itemize}[nosep]
  \item Well~A (good generalization): $H_A$, $\Sigma_A$, with
        $\Sigma_A / H_A = 1$.
  \item Well~B (bad generalization): $H_B$, $\Sigma_B$, with
        $\Sigma_B / H_B = 1$.
\end{itemize}
The loss values $L_A = L_B$ and the landscape is symmetric.  Despite well~A
having lower generalization gap (due to different higher-order structure or
sample size dependence not captured by $\Sigma/H$ alone), the noise ratios
are equal.

\begin{proof}
By \Cref{prop:fpe}, the stationary density ratio at the two minima is
\[
  \frac{p(\theta_A)}{p(\theta_B)}
  = \frac{\Sigma_B}{\Sigma_A}
    \cdot\exp\!\Bigl(-\frac{2}{\tau}\int_{\theta_B}^{\theta_A}
      \frac{L'(s)}{\Sigma(s)}\dif s\Bigr).
\]
In the symmetric case, the integral $\int_{\theta_B}^{\theta_A}
L'(s)/\Sigma(s)\dif s = 0$ by antisymmetry.  Hence
\[
  \frac{p(\theta_A)}{p(\theta_B)} = \frac{\Sigma_B}{\Sigma_A}.
\]
Since $\Sigma_A/H_A = \Sigma_B/H_B$ with equal Hessians in the symmetric
construction (one can set $H_A = H_B$), we have $\Sigma_A = \Sigma_B$, giving
$p(\theta_A)/p(\theta_B) = 1$.

More generally, even without $H_A = H_B$, the condition $\Sigma_A/H_A =
\Sigma_B/H_B$ means the effective potential curvatures
$\Veff''(\theta_k) \approx H_k/\Sigma_k$ are equal at both wells.  In a
Laplace approximation of the stationary distribution near each well,
\[
  \pi(k) \propto \frac{1}{\Sigma_k}\sqrt{\frac{2\pi\tau}{2H_k/\Sigma_k}}
  = \frac{1}{\Sigma_k}\sqrt{\frac{\pi\tau\Sigma_k}{H_k}}
  = \frac{1}{\sqrt{\Sigma_k}}\sqrt{\frac{\pi\tau}{H_k}},
\]
and with $\Sigma_A/H_A = \Sigma_B/H_B$ and the symmetric construction
$L_A = L_B$:
\[
  \frac{\pi(A)}{\pi(B)}
  = \frac{\sqrt{\Sigma_B/H_B}}{\sqrt{\Sigma_A/H_A}} = 1.
\]
Thus $\pi(A) \approx \pi(B) \approx 1/2$.  SGD selects the good solution
with probability no better than a fair coin.
\end{proof}

\paragraph{Which condition is violated.}
Condition~N3 (the noise ratios $\Sigma/H$ must differ across minima) is
violated.  When all minima share the same ratio, the noise-driven selection
mechanism provides no bias toward good solutions.

% ── 3.4 F4 ───────────────────────────────────────────────
\subsection{F4: Noise-Generalization Reversal (Main Result)}
\label{sec:f4}

\begin{theorem}[Noise-generalization reversal]
\label{thm:f4}
There exists a one-dimensional loss landscape with two local minima---one with
low generalization gap (good) and one with high generalization gap (bad)---such
that SGD converges to the bad solution with overwhelming probability.
Specifically, the ratio of escape times satisfies
$\tau_A / \tau_B \sim \exp(-19/\eta) \to 0$ as $\eta\to 0$.
\end{theorem}

\paragraph{Construction.}
Define a double-well loss with equal loss values and equal curvatures but
different noise variances:
\begin{itemize}[nosep]
  \item Well~A (good generalization): $L_A = 0.1$, $H_A = 2$,
        $\Sigma_A = 20$.
  \item Well~B (bad generalization): $L_B = 0.1$, $H_B = 2$,
        $\Sigma_B = 1$.
\end{itemize}
The generalization gaps (\Cref{prop:gap}) are
\[
  \mathrm{Gap}_A = \frac{\Sigma_A}{H_A(n-1)} = \frac{20}{2(n-1)} = \frac{10}{n-1},
  \qquad
  \mathrm{Gap}_B = \frac{\Sigma_B}{H_B(n-1)} = \frac{1}{2(n-1)} = \frac{0.5}{n-1}.
\]
So well~A (the good solution in terms of diversity and feature usage) has a
\emph{larger} generalization gap as measured by $\Sigma/H$---this is
deliberately reversed to construct the counterexample.

We pause to clarify a subtle but critical point.  The LOO proxy
$\mathrm{Gap} \approx \Sigma/(H(n-1))$ from \Cref{prop:gap} predicts that
well~A has a \emph{larger} generalization gap than well~B.  This is precisely
the point of the counterexample: the LOO proxy is an approximation derived
under regularity assumptions, and it measures sensitivity to single-sample
removal, not actual population risk.

The construction posits a data-generating process where well~A corresponds
to a solution that leverages diverse features---different training examples
activate different model components, producing diverse per-sample gradients
(hence high $\Sigma_A = 20$) but low population risk because the diverse
features capture the true data structure.  Well~B corresponds to a memorizing
solution that fits a uniform low-frequency pattern, producing nearly
identical per-sample gradients (hence low $\Sigma_B = 1$) but high
population risk because the uniform pattern does not capture the true
structure.  The LOO proxy $\Sigma/(H(n-1))$ fails here because high
per-sample gradient diversity does not imply instability to sample
removal---it reflects feature diversity, not fragility.

This is the essence of condition~N4: the noise variance $\Sigma$ tracks
per-sample sensitivity, not generalization \emph{per se}.  When the
correlation between $\Sigma$ and actual generalization gap is negative,
SGD's low-$\Sigma$ preference systematically selects bad solutions.

\begin{proof}
We compute the effective potential barriers and Kramers escape times for each
well.

\medskip\noindent\textbf{Step~1: Effective potential near each well.}
From~\eqref{eq:veff}, the effective potential near a minimum $\theta_k$
(where $L'(\theta_k) = 0$) has the local expansion
\[
  \Veff(\theta)
  \approx \Veff(\theta_k)
    + \frac{1}{2}\Veff''(\theta_k)(\theta - \theta_k)^2 + \cdots
\]
The second derivative of $\Veff$ at a minimum is
\[
  \Veff''(\theta_k)
  = \frac{H_k}{\Sigma_k}
    + \frac{\tau}{2}\frac{\Sigma''_k\Sigma_k - (\Sigma'_k)^2}{\Sigma_k^2}.
\]
In the leading-order ($\tau\to 0$) analysis, the dominant term is $H_k/\Sigma_k$.

For wells~A and~B:
\[
  \Veff''(\theta_A) \approx \frac{H_A}{\Sigma_A} = \frac{2}{20} = 0.1,
  \qquad
  \Veff''(\theta_B) \approx \frac{H_B}{\Sigma_B} = \frac{2}{1} = 2.
\]

\medskip\noindent\textbf{Step~2: Effective barrier heights.}
Let $\theta_s$ denote the saddle point between the two wells, with
$L(\theta_s) = L_A + \Delta V = L_B + \Delta V$ for some barrier height
$\Delta V > 0$ in the original loss.

The effective potential barrier from well~$k$ to the saddle is, to leading
order,
\[
  \Delta U_k
  = \Veff(\theta_s) - \Veff(\theta_k)
  \approx \frac{\Delta V}{\Sigma_k},
\]
where we have used the fact that the integral
$\int_{\theta_k}^{\theta_s} L'(s)/\Sigma(s)\dif s$ is dominated by the
behavior near well~$k$ where $\Sigma\approx\Sigma_k$.

We choose a concrete barrier height $\Delta V = 10$ in the original loss:
\[
  \Delta U_A = \frac{\Delta V}{\Sigma_A} = \frac{10}{20} = 0.5,
  \qquad
  \Delta U_B = \frac{\Delta V}{\Sigma_B} = \frac{10}{1} = 10.
\]
The ratio of effective barriers is
\[
  \frac{\Delta U_A}{\Delta U_B}
  = \frac{\Sigma_B}{\Sigma_A}
  = \frac{1}{20}.
\]
The good well~A has a barrier 20 times \emph{lower} in the effective potential
than the bad well~B.

\medskip\noindent\textbf{Step~3: Kramers escape times.}
The Kramers escape time from well~$k$ is \citep{kramers1940brownian}
\[
  \tau_k \sim \frac{1}{\eta}\exp\!\Bigl(\frac{2\,\Delta U_k}{\eta}\Bigr).
\]
Taking the ratio with $\Delta U_A = 0.5$ and $\Delta U_B = 10$:
\begin{equation}\label{eq:f4-ratio}
  \frac{\tau_A}{\tau_B}
  \sim \exp\!\Bigl(\frac{2(\Delta U_A - \Delta U_B)}{\eta}\Bigr)
  = \exp\!\Bigl(\frac{2(0.5 - 10)}{\eta}\Bigr)
  = \exp\!\Bigl(-\frac{19}{\eta}\Bigr).
\end{equation}

\medskip\noindent\textbf{Step~4: Stationary distribution.}
In the two-well approximation, the stationary occupation probabilities satisfy
\[
  \frac{\pi(B)}{\pi(A)}
  = \frac{\tau_B}{\tau_A}
  \sim \exp\!\Bigl(\frac{19}{\eta}\Bigr) \to \infty
  \quad\text{as } \eta\to 0.
\]
Therefore $\pi(B) \to 1$: SGD converges to the bad solution~B with
overwhelming probability.

\medskip\noindent\textbf{Step~5: Verification of the reversal.}
Well~B has the lower gradient noise variance ($\Sigma_B = 1 < 20 = \Sigma_A$),
so SGD's noise-driven preference favors~B.  But well~B has worse
generalization.  The stationary distribution
$p(\theta_B)/p(\theta_A) = (\Sigma_A/\Sigma_B)\cdot(\cdots) = 20\cdot(\cdots)$,
where the exponential factor further amplifies the preference for~B.
This is the noise-generalization reversal.
\end{proof}

\paragraph{Why this happens.}
The reversal has a natural interpretation.  The good solution~A achieves low
test error by leveraging diverse features in the data.  Different training
examples activate different features, leading to diverse per-sample gradients
and hence high $\Sigma_A$.  The bad solution~B memorizes a uniform pattern
(e.g., a low-frequency component or a label-correlated shortcut), producing
nearly identical per-sample gradients and hence low~$\Sigma_B$.

SGD's noise-driven selection mechanism~\eqref{eq:fpe-steady} favors regions
of low $\Sigma$, which is precisely the memorizing solution.  This is the
fundamental limit of noise-driven implicit regularization: the same mechanism
that in typical settings provides beneficial regularization can, when the
correlation between noise variance and generalization is reversed,
systematically select bad solutions.

\paragraph{Which condition is violated.}
Condition~N4 ($\mathrm{Corr}(\Sigma, \mathrm{Gap}) > 0$) is violated.
In this construction, high~$\Sigma$ corresponds to \emph{low} generalization
gap (good) and low~$\Sigma$ corresponds to high generalization gap (bad),
reversing the alignment that implicit regularization requires.

% ══════════════════════════════════════════════════════════
\section{Necessary Conditions Framework}
\label{sec:necessary}

The four failure modes of \Cref{sec:failure} directly yield four necessary
conditions for SGD's implicit regularization to succeed.

\begin{theorem}[Four necessary conditions]
\label{thm:necessary}
For SGD to converge to a good solution (in the sense of low generalization gap)
with high probability, the following conditions are all necessary:
\begin{enumerate}[label=\textup{\textbf{(N\arabic*)}},leftmargin=3em]
  \item \textbf{Sufficient mixing:}
        $\eta T \gg \tau_{\mathrm{mix}}$, where $\tau_{\mathrm{mix}}$ is the
        mixing time of the SGD Markov chain.
  \item \textbf{Discrete stability:}
        $\eta < 2/\lambda_{\max}$, where $\lambda_{\max}$ is the largest
        per-sample Hessian eigenvalue.
  \item \textbf{Selection power:}
        The noise-to-curvature ratios $\Sigma(\theta^*)/H(\theta^*)$ differ
        across local minima.
  \item \textbf{Correct alignment:}
        $\mathrm{Corr}(\Sigma, \mathrm{Gap}) > 0$, i.e., minima with higher
        gradient noise variance also have higher generalization gap.
\end{enumerate}
\end{theorem}

\begin{proof}
Each condition is necessary because violating it leads to a demonstrated
failure:
\begin{itemize}[nosep]
  \item Violating~(N1) leads to failure mode~F1 (\Cref{prop:f1}): SGD is
        trapped in the initialization basin and cannot reach the good solution,
        regardless of the stationary distribution's preference.
  \item Violating~(N2) leads to failure mode~F2 (\Cref{prop:f2}): SGD diverges
        due to a positive Lyapunov exponent, preventing convergence to any
        solution.
  \item Violating~(N3) leads to failure mode~F3 (\Cref{prop:f3}): the
        stationary distribution assigns equal weight to good and bad minima,
        so SGD provides no selective advantage.
  \item Violating~(N4) leads to failure mode~F4 (\Cref{thm:f4}): the
        noise-driven selection systematically favors the bad solution.
\end{itemize}
In each case, the failure is demonstrated by an explicit construction, so
the condition is necessary for success across all problem instances.
\end{proof}

\begin{remark}[The conjunction is not proven sufficient]
\label{rem:not-sufficient}
While conditions~(N1)--(N4) are each necessary, we do not claim that their
conjunction is sufficient for SGD to find good solutions.  Additional
conditions---related to the landscape geometry, initialization distribution,
or higher-order properties of the noise structure---may also be required.
We leave the question of identifying a minimal set of sufficient conditions
as an important open problem.
\end{remark}

\begin{remark}[Three-level classification of conditions]
\label{rem:three-level}
The conditions can be organized into a three-level hierarchy:
\begin{enumerate}[nosep]
  \item \textbf{Necessary:} The global minimum of the effective potential
        $\Veff$ lies in the set of good solutions.  If this fails, the
        stationary distribution concentrates on a bad solution (as
        $\tau\to 0$).
  \item \textbf{Sufficient:} The global minimum of $\Veff$ is unique in the
        good set, with spectral gap $\Delta > 0$, and $\tau \ll \Delta$.
        Then $\Prob(\text{good solution}) \ge 1 - (K-1)\exp(-2\Delta/\tau)$.
  \item \textbf{Approximately sufficient:} All minima of $\Veff$ with values
        within $\delta$ of the global minimum are approximately good (gap
        $\le \varepsilon + \varepsilon'$).  Then
        $\Prob(\mathrm{Gap}\le\varepsilon+\varepsilon') \ge 1-\gamma$, where
        $\gamma \sim K e^{-2\delta/\tau}$.
\end{enumerate}
Conditions~(N1)--(N4) are at level~1.  Closing the gap to level~2 or~3
requires quantitative control over the landscape that goes beyond our
current results.
\end{remark}

% ══════════════════════════════════════════════════════════
\section{Reparameterization Invariance}
\label{sec:reparam}

A natural concern is whether the failure mode~F4 is an artifact of a
particular parameterization.  We show that it is not, because the key
quantity $\Sigma/H$ is reparameterization-invariant.

\begin{proposition}[Reparameterization invariance of $\Sigma/H$]
\label{prop:reparam}
Let $\theta = g(\varphi)$ be a smooth bijection (reparameterization), and
let $\varphi^* = g^{-1}(\theta^*)$.  Define the reparameterized Hessian
$\tilde{H} = L''(g(\varphi^*))\cdot [g'(\varphi^*)]^2$ and the
reparameterized noise variance
$\tilde{\Sigma} = \frac{1}{n}\sum_i [\ell_i'(g(\varphi^*))]^2\cdot [g'(\varphi^*)]^2$.
Then
\[
  \frac{\tilde{\Sigma}}{\tilde{H}} = \frac{\Sigma(\theta^*)}{H(\theta^*)}.
\]
\end{proposition}

\begin{proof}
Under the reparameterization $\theta = g(\varphi)$, the chain rule gives
\[
  \frac{\partial \ell_i}{\partial\varphi}\bigg|_{\varphi^*}
  = \ell_i'(g(\varphi^*))\cdot g'(\varphi^*)
  = \ell_i'(\theta^*)\cdot g'(\varphi^*).
\]
Therefore the reparameterized noise variance is
\[
  \tilde{\Sigma}
  = \frac{1}{n}\sum_{i=1}^{n}
    \Bigl[\ell_i'(\theta^*)\cdot g'(\varphi^*)\Bigr]^2
  = \Sigma(\theta^*)\cdot \bigl[g'(\varphi^*)\bigr]^2.
\]
Similarly, the reparameterized Hessian of $L$ at $\varphi^*$ (using
$L'(\theta^*)=0$, so the first-derivative term in the chain rule for
second derivatives vanishes) is
\[
  \tilde{H}
  = L''(\theta^*)\cdot\bigl[g'(\varphi^*)\bigr]^2
  = H(\theta^*)\cdot\bigl[g'(\varphi^*)\bigr]^2.
\]
The ratio is
\[
  \frac{\tilde{\Sigma}}{\tilde{H}}
  = \frac{\Sigma(\theta^*)\cdot[g'(\varphi^*)]^2}{H(\theta^*)\cdot[g'(\varphi^*)]^2}
  = \frac{\Sigma(\theta^*)}{H(\theta^*)}.
  \qedhere
\]
\end{proof}

\begin{corollary}
\label{cor:gap-intrinsic}
The generalization gap proxy $\mathrm{Gap}(\theta^*) \approx \Sigma(\theta^*)/(H(\theta^*)(n-1))$ from \Cref{prop:gap} is a reparameterization-invariant quantity.
\end{corollary}

\begin{proof}
Immediate from \Cref{prop:reparam}: $\tilde{\Sigma}/\tilde{H} = \Sigma/H$,
so $\tilde{\Sigma}/(\tilde{H}(n-1)) = \Sigma/(H(n-1))$.
\end{proof}

\paragraph{Implication for F4.}
The noise-generalization reversal in \Cref{thm:f4} is not an artifact of a
particular coordinate system.  Since both $\Sigma/H$ and the effective
potential $\Veff$ are built from reparameterization-invariant quantities, the
conflict between SGD's noise preference and generalization persists under any
smooth reparameterization.  This underscores that F4 is a fundamental
limitation, not a coordinate-dependent pathology.

% ══════════════════════════════════════════════════════════
\section{Discussion}
\label{sec:discussion}

% ── 6.1 ──────────────────────────────────────────────────
\subsection{Implications for Practice}
\label{sec:practice}

Each failure mode suggests a practical intervention:
\begin{itemize}[nosep]
  \item \textbf{F1 (insufficient mixing)}: Use learning rate warm-up or
        cyclical learning rate schedules to help SGD escape bad basins early
        in training.
  \item \textbf{F2 (discrete instability)}: Do not use learning rates that
        are too large.  The stability threshold $\eta < 2/\lambda_{\max}$
        provides a principled upper bound.
  \item \textbf{F3 (degenerate selection)}: Data augmentation and diverse
        training data can break symmetry in $\Sigma/H$ ratios across minima,
        restoring SGD's selective power.
  \item \textbf{F4 (noise-generalization reversal)}: When the correlation
        between gradient noise and generalization is negative---as may occur
        with small datasets, shortcut features, or distribution shift---explicit
        regularization (weight decay, dropout, early stopping) is essential.
        SGD's implicit regularization alone is insufficient.
\end{itemize}

% ── 6.2 ──────────────────────────────────────────────────
\subsection{Connection to the Interpolation Regime}
\label{sec:interpolation}

In the modern over-parameterized interpolation regime, all training losses
reach zero: $\ell_i(\theta^*)=0$ for all~$i$.  This implies
$\nabla\ell_i(\theta^*)=0$ and hence $\Sigma(\theta^*)=0$ at every
interpolating solution.  The continuous-time SDE framework becomes degenerate,
and the noise-driven selection mechanism analyzed here does not directly apply.

Instead, implicit regularization in the interpolation regime operates through
a different mechanism: discrete stability.  Solutions with
$\eta\,h_{\max} > 2$ are destabilized, while those with $\eta\,h_{\max} < 2$
remain stable \citep{wu2023implicit}.  This selects for solutions with low
Jacobian complexity $\mathcal{C}(\theta^*) = \|J\|_F/\sqrt{n}$.

Our failure mode~F2 is relevant in this regime: choosing $\eta$ too large
may destabilize all solutions.  Failure modes F3 and F4, as stated, pertain
primarily to the non-interpolation regime where $\Sigma > 0$.  However, the
conceptual lesson of~F4---that the metric optimized by SGD's dynamics may
not align with generalization---extends to the interpolation regime as well:
low Jacobian complexity is not a universal proxy for generalization.

% ── 6.3 ──────────────────────────────────────────────────
\subsection{Relation to Prior Work}
\label{sec:prior}

\paragraph{Shape matters \citep{li2021happens}.}
Li et al.\ demonstrate that SGD's noise structure (not just its magnitude)
influences which minima are selected.  Our work is complementary: we identify
four settings where even the correct noise structure fails to produce good
solutions.

\paragraph{Sharp minima can generalize \citep{dinh2017sharp}.}
Dinh et al.\ show that reparameterization can make any minimum arbitrarily
sharp without affecting generalization, challenging flatness-based explanations.
Our \Cref{prop:reparam} resonates with this observation: $H$ alone is not
invariant, but $\Sigma/H$ is.  However, F4 shows that even the
reparameterization-invariant quantity $\Sigma/H$ can fail as a generalization
proxy.

\paragraph{Implicit regularization in linear models \citep{arora2019implicit}.}
Arora et al.\ characterize implicit regularization for deep linear networks.
Their results show that architecture (depth, width) shapes the implicit bias.
Our F4 highlights that even with fixed architecture, the data distribution
can induce a noise-generalization reversal.

\paragraph{Entropy-SGD \citep{chaudhari2018stochastic}.}
Chaudhari and Soatto connect SGD to an entropy-regularized objective,
establishing a link between noise and flatness-seeking behavior.  Our
\Cref{prop:fpe} uses the same FPE framework; F4 demonstrates a failure
of the resulting selection when the noise-generalization correlation is
negative.

% ── 6.4 ──────────────────────────────────────────────────
\subsection{Open Questions}
\label{sec:open}

\begin{enumerate}[nosep]
  \item \textbf{Sufficiency.}  Are conditions (N1)--(N4) jointly sufficient?
        We conjecture they are not, and that additional landscape conditions
        are needed.  Identifying a minimal sufficient set is a key open problem.
  \item \textbf{Multi-dimensional effective potential.}  In dimension $d>1$,
        the detailed balance condition may fail, and the steady state involves
        a Freidlin--Wentzell quasipotential rather than an explicit formula.
        Extending our results to this setting is technically demanding.
  \item \textbf{Prevalence of F4.}  How often does the noise-generalization
        reversal occur in practical deep learning?  Characterizing the
        data-architecture combinations that lead to $\mathrm{Corr}(\Sigma,
        \mathrm{Gap}) < 0$ is an empirical and theoretical challenge.
  \item \textbf{Global trajectory analysis.}  Our results are local (near
        minima).  Understanding SGD's global trajectory from random
        initialization to a basin of attraction requires different tools.
  \item \textbf{Discrete-continuous unification.}  The non-interpolation and
        interpolation regimes use fundamentally different mathematical
        frameworks.  A unified theory handling the $\Sigma\to 0$ limit
        remains open.
\end{enumerate}

% ══════════════════════════════════════════════════════════
\section{Conclusion}
\label{sec:conclusion}

We have identified four failure modes of SGD's implicit regularization, each
demonstrated by an explicit construction with quantitative analysis.  Failure
modes F1--F3 concern the dynamics (mixing time, stability, and selection
power), while F4---our main result---reveals a fundamental limitation: SGD's
noise-driven preference for low gradient variance can systematically select
solutions that generalize poorly when good solutions use diverse features
(high $\Sigma$) and bad solutions memorize (low $\Sigma$).

From these failure modes we derived four necessary conditions (N1--N4) that
must hold simultaneously for implicit regularization to succeed.  The
conditions span three distinct aspects: the optimization dynamics (N1, N2),
the noise structure (N3), and the data-dependent alignment between noise and
generalization (N4).  Their conjunction is not shown to be sufficient, and
closing this gap remains an important open problem.

Our results carry a practical message: in settings where condition~N4 may be
violated---small datasets, distribution shift, shortcut features---relying
solely on SGD's implicit regularization is insufficient, and explicit
regularization is needed.


% ══════════════════════════════════════════════════════════
\begin{thebibliography}{10}

\bibitem[Arora et~al.(2019)]{arora2019implicit}
S.~Arora, N.~Cohen, W.~Hu, and Y.~Luo.
\newblock Implicit regularization in deep matrix factorization.
\newblock In \emph{Advances in Neural Information Processing Systems}, 2019.

\bibitem[Chaudhari and Soatto(2018)]{chaudhari2018stochastic}
P.~Chaudhari and S.~Soatto.
\newblock Stochastic gradient descent performs variational inference, converges
  to limit cycles for deep networks.
\newblock In \emph{International Conference on Learning Representations}, 2018.

\bibitem[Dinh et~al.(2017)]{dinh2017sharp}
L.~Dinh, R.~Pascanu, S.~Bengio, and Y.~Bengio.
\newblock Sharp minima can generalize for deep nets.
\newblock In \emph{International Conference on Machine Learning}, 2017.

\bibitem[Horsthemke and Lefever(1984)]{horsthemke1984noise}
W.~Horsthemke and R.~Lefever.
\newblock \emph{Noise-Induced Transitions: Theory and Applications in Physics,
  Chemistry, and Biology}.
\newblock Springer-Verlag, 1984.

\bibitem[Kramers(1940)]{kramers1940brownian}
H.~A. Kramers.
\newblock Brownian motion in a field of force and the diffusion model of
  chemical reactions.
\newblock \emph{Physica}, 7(4):284--304, 1940.

\bibitem[Li et~al.(2021)]{li2021happens}
Z.~Li, T.~Wang, and S.~Arora.
\newblock What happens after {SGD} reaches zero loss? {A} mathematical
  framework.
\newblock In \emph{International Conference on Learning Representations}, 2021.

\bibitem[Smith and Le(2018)]{smith2018dont}
S.~L. Smith and Q.~V. Le.
\newblock A {B}ayesian perspective on generalization and stochastic gradient
  descent.
\newblock In \emph{International Conference on Learning Representations}, 2018.

\bibitem[Wu and Su(2023)]{wu2023implicit}
L.~Wu and W.~J. Su.
\newblock The implicit regularization of dynamical stability in stochastic
  gradient descent.
\newblock In \emph{Proceedings of the 40th International Conference on Machine
  Learning (ICML)}, volume 202 of \emph{PMLR}, pp.\ 37656--37684, 2023.

\end{thebibliography}

\end{document}
